\section{Introduction}\label{sec:introduction}

Matrix multiplication asks for the bilinear map \((A,B)\mapsto AB\),
where the two input matrices have compatible sizes.  Its arithmetic
complexity measures the number of field operations needed to compute
all entries of the product.  Except in the field-transfer results
of Section~\ref{sec:characteristic-transfer}, the field is \(\C\), and
the constants in an algorithm may belong to \(\C\).
For \(0\le k\le1\), let
\[
 w(k)=\omega(1,k,1)
\]
be the infimum of exponents for multiplying an
\(n\times\lceil n^k\rceil\) matrix by a
\(\lceil n^k\rceil\times n\) matrix.
The bound \(w(k)\le\tau\) means that for every \(\varepsilon>0\)
these products can be computed using \(O(n^{\tau+\varepsilon})\)
arithmetic operations.
The square exponent is \(\omega=w(1)\), and the dual exponent is
\[
 \alpha=\sup\{k\in[0,1]:w(k)=2\}.
\]
The output size gives \(w(k)\ge2\).  Thus \(\alpha\) measures how far
the inner dimension can grow while multiplication retains the exponent
of its output size.  Square and rectangular multiplication are basic
operations in algebraic algorithms; controlling the rectangular shape
can matter even when a square bound is already available.

\subsection{Context and main result}

Strassen's rank-seven algorithm for \(2\times2\) multiplication gave
\(\omega\le\log_2 7\), establishing subcubic arithmetic complexity
\cite{strassen}.  Approximate bilinear algorithms and their conversion
to exact algorithms \cite{bini}, Sch\"onhage's asymptotic sum inequality
\cite{schonhage}, and Strassen's laser method \cite{strassen-laser}
made tensor powers and degenerations central to subsequent advances.
Coppersmith and Winograd combined low-border-rank tensors with
arithmetic-progression constructions \cite{cw}.  Their extraction
method restricts powers of a small tensor to direct sums of matrix
products, whose dimensions enter the asymptotic sum inequality;
their tensor-square analysis gave \(\omega<2.375477\)
\cite[Section~8]{cw}.
Higher powers and more detailed analyses of these tensors
led to the improvements of Stothers, Vassilevska Williams, and Le Gall
\cite{stothers,williams-powers,legall-powers}.

Rectangular multiplication requires separate control of the three
matrix dimensions.  Early work of Coppersmith and of Lotti and Romani
developed this direction \cite{coppersmith1982,lottiromani1983}; later
constructions of Coppersmith, Le Gall, and Le Gall and Urrutia improved
the available rectangular and dual bounds
\cite{coppersmith-rect,legall-rect,legall-urrutia}.
Recent advances include the refined laser method of Alman and
Vassilevska Williams \cite{alman-williams}, the asymmetric hashing of
Duan, Wu, and Zhou \cite{duan-wu-zhou}, and the square and rectangular
analyses of Vassilevska Williams, Xu, Xu, and Zhou \cite{vxxz} and
Alman, Duan, Vassilevska Williams, Xu, Xu, and Zhou \cite{advxxz}.
The four-author work establishes \(\alpha\ge0.321334\); these two
papers also give detailed tables of rectangular estimates.
Dupont et al.\ subsequently optimized the asymmetric framework at
a larger scale to obtain \(\omega<2.371177\)
\cite{dupont}, a bound incorporated into the updated table of
\cite{advxxz}.

The common extraction problem is to separate many candidate products
that share variables, while accounting for the cost of every operation
used to separate them.  We adjoin auxiliary group tensors that both
resolve shared labels and supply useful matrix-product factors.  Their
full rank cost enters the resulting inequalities.  This construction
gives the following bounds.

\begin{theorem}\label{thm:main}\label{thm:square}\label{thm:rectangular}
Over \(\C\), the arithmetic matrix-multiplication exponents satisfy
\[
 \omega<\frac{1129}{500}=2.258,\qquad
 \alpha>0.465,\qquad w(0.709)<2.092.
\]
\end{theorem}

Corollary~\ref{cor:characteristic-transfer} transfers the strict square
and \(k=0.709\) inequalities to every field of characteristic zero or
outside one finite set of positive characteristics, in the same
arithmetic-operation model.  It selects two fixed exact rank schemes,
realizes their coefficients over one number field, and specializes
their polynomial identities.  The exceptional primes are not computed.
Corollary~\ref{cor:characteristic-zero-dual} also gives the same
dual-exponent bound over every field of characteristic zero.
No dual-exponent transfer in positive characteristic is asserted.

The square bound also lowers the arithmetic exponent of basic
linear-algebra tasks.  Given \(A\in\C^{n\times n}\) and
\(b\in\C^n\), one can compute \(\det A\) and, when \(A\) is
nonsingular, \(A^{-1}\) and the solution of \(Ax=b\), using
\(O(n^c)\) field operations for some fixed \(c<2.258\), with exact
zero tests for pivot selection.  Choose \(\omega<c<2.258\) and apply
the fast elimination reductions of Strassen \cite{strassen} and
Ibarra, Moran, and Hui
\cite[Theorem~2.1 and Corollary~2.1]{ibarra-moran-hui}.

Two applications of the same separation lemma prove the theorem.
The square bound uses finite tensor constructions, a differential
comparison for their attainable values, and an exact numerical
certificate.  The dual and rectangular bounds come from explicit
conditional-label trees and rational entropy certificates.
We also optimize all leaf probability laws in a specified finite
family of such trees, obtaining a bound for every aspect ratio in
\([0,1]\).  The square and rectangular construction families, and
the precise way a square estimate augments the latter, are described
below.

\subsection{Separating labels with an auxiliary tensor}

A trilinear tensor is a sum of monomials \(xyz\), with variables
drawn from three finite coordinate sets \(X,Y,Z\).  Its rank is the
least number of products of three linear forms needed to express it.
A restriction substitutes linear forms separately on the three
sides.  A degeneration allows parameter-dependent substitutions and
takes a leading coefficient.  Border rank is the least rank of
tensors approaching the given tensor.  These operations are local: a
coordinate substitution cannot depend on information visible only
on another side.

Suppose that the supported monomials carry \(K\) labels, and that
each of two sides can independently determine the label from its own
coordinate.  The third side may still share variables between labels.
The pairwise separation lemma, \cref{lem:group}, tensors the input
with a finite abelian group-addition tensor of rank \(K^{1+o(1)}\).
Local projections then give a direct sum over the labels, with a
size-\(K\) dot-product factor on the two original readers in every
summand.  The label is now visible on all three sides.  Dot-product
factors on the three possible reader pairs combine into a matrix
multiplication tensor.

The construction combines Fourier diagonalization of the group tensor
with tags of equal Euclidean norm and a separate constant-inner-product
slice for each tag.  Its geometry adapts Behrend's construction
\cite{behrend}, following the progression-free constructions of Salem
and Spencer \cite{salem-spencer}.  The use of group algebras also has
an established place in the approach of Cohn and Umans
\cite{cohn-umans}; the local separation identity needed here is proved
in full.  It preserves every original coordinate and coefficient,
including additional tensor indices correlated across occurrences.
Consequently the same identity can be applied inside a previously
restricted tensor.  A separate support argument ensures that every
summand counted in an application is nonzero.

Successive separations are organized by exact conditional counts.
Once a preceding label is direct on all three sides, each side can
locate the positions on which the next operation acts.  Fixing the
number of each possible continuation makes the next auxiliary tensor
and the output dimensions common to all branches.  The multinomial
counts and entropy chain rule are the standard method-of-types
tools \cite{shannon1948,csiszar1998}; locality and the uniform tensor
accounting are proved alongside them here.

\subsection{The square bound: finite constructions and comparison}

The square argument starts from the six-term Coppersmith--Winograd
tensor
\[
 T=\sum_{\substack{i+j+k=2\\0\le i,j,k\le2}}x_i y_j z_k,
\]
which has border rank three \cite[Section~7, equation~(10)]{cw}.
A five-term degeneration of \(T\) has three groups, each detectable
on one designated side.  A local degeneration of three copies retains
\(75\) scalar components.  Separating their successive labels gives
the elementary estimate
\(\omega\le9\log3/\log75<2.290108\)
in \cref{thm:elementary}.  This example displays the auxiliary-rank
accounting before the larger construction is introduced.

For a proposed exponent \(\beta\), suppose a finite construction
from \(n\) copies of \(T\) and an auxiliary tensor of border rank
at most \(R\) produces \(L\) independent matrix products of common
dimensions \(a\times b\) by \(b\times c\).  Its normalized score is
\[
 \frac{\log L-\log R+(\beta/3)\log(abc)}{n}.
\]
Sch\"onhage's asymptotic sum inequality implies that a score strictly
greater than \(\log3\) proves \(\omega<\beta\).
Thus the square argument must produce a strict score gain while
retaining the entire auxiliary charge.

We organize the finite constructions by their source distribution.
There are six collections of original tensor occurrences, called
\emph{orientation banks}, one for each permutation of the three
sides.  The distribution records the same exact frequencies in the
original coordinates of each bank.  This keeps a physical permutation
of a tensor distinct from a change of its recorded law.  The square
sections use \(\alpha\) for this vector-valued source law; the dual
exponent in \cref{thm:main} is a scalar.  A construction
preserves arbitrary unresolved tensors attached to its original words,
and its output labels identify distinct surviving words.
Taking the closure of the finite law--score pairs gives a continuous
concave attainable-value function.  Conditional pooling and a new
application of the group separator to each longer pool provide the
composition properties needed in the comparison.

The key analytic step uses restrictions on sums of the local grades
\(i,j,k\).  Each side can read its own sum.  Windows placed at
progression-free target indices force the three locally read targets
to agree, creating another direct label.  The progression-free
construction and the matching argument have the ancestry of Behrend
and Coppersmith--Winograd \cite{behrend,cw}.
To exploit the new labels, a finite-horizon control steers every word
allowed by its controls into the assigned terminal window.  This
every-word guarantee is what permits all the resulting tensor
summands to be counted.  Exact rational prefix populations then
translate the finite control tree into local tensor maps with common
auxiliary supplies.

This construction forces a differential inequality on any smooth
function touching the attainable value from below.  The covariance
of the local grades determines the weighted combination of second
derivatives in this inequality.  Smooth touching
tests and strict comparison follow the viscosity-solution methodology
of Crandall and Lions and of Crandall, Ishii, and Lions
\cite{crandalllions1983,crandallishiilions1992}; the required
inequality and comparison are proved directly from the finite tensor
constructions.  A rational polynomial chart and a polynomial barrier
then give the strict score gain at an interior source distribution.
Boundary constructions supply the initial lower estimates, while
explicit residual, interpolation, and rounding bounds turn the finite
integer checks into inequalities on the entire comparison domain.

Finally, a strict gain in the closed attainable set selects one
actual finite construction with score greater than \(\log3\).
Its law may be close to the chart's central law.  All control horizons,
rational approximations, group sizes, and bank sizes are fixed before
the asymptotic sum inequality is applied.

\subsection{Rectangular bounds and finite-family optimization}

The rectangular argument tracks the three matrix dimensions separately.
A readable label tree partitions a tensor's support until each leaf
is one monomial; at every node, conditional on the preceding labels,
the next child is independently readable on two specified sides.
Fix a rational probability law on its leaves and apply the tree to
an admissible number of copies with those exact leaf frequencies;
arbitrary leaf laws follow by approximation.  The number of child words
at a node is a multinomial coefficient.  The product of these
coefficients is exactly the number of final direct summands, so the
leading auxiliary rank charge is canceled by this multiplicity in
the equal-summand inequality.  The three remaining logarithmic
dimensions are the conditional entropies accumulated on the three
reader pairs.  This is \cref{thm:entropy}.

Explicit polynomial degenerations giving border-rank bounds two and
three supply the initial trees.  Their supports have three colors, with a
designated side able to detect its own color.  A \emph{completion}
takes several copies, chooses a center color, and retains color words
containing at most one distinct noncenter color.  Local color tests
implement this leading-term degeneration and preserve the conditional trees.
Deleting one color at the end gives a fully readable tree.
The family \(\calF_7\) consists of both bases, at most seven
completions of sizes two or three, all center choices, all pairs of
retained colors, and all choices of a distinguished side.

For each member of this finite family, a three-state recurrence
maximizes every nonnegative weighted sum of its incident and
nonincident conditional entropies over all leaf laws.  Its support
function therefore describes the closed downward convex hull of the
normalized rate pairs exactly.  \Cref{thm:shape-optimization} turns
this description into a maximization over one real parameter for
each aspect ratio.  An available square estimate contributes one
additional rate point.  In particular, the square part of
\cref{thm:main} supplies the point for \(p=2.258\).
This finite-family optimization is distinct from the square
attainable-value construction with orientation banks and count windows.

Two explicit laws give the dual and rectangular conclusions.
For the dual exponent, an exact conditional-independence identity
preserves all the entropy of a uniform coordinate word on its incident
reader pairs.  This equality certifies exponent two, and the remaining
entropy certifies the inner dimension.  A small fiber histogram gives
the rectangular point.  Rational enclosures establish their strict
numerical margins.  The same finite family also gives an independent
square bound below \(2.267\) by an integer leaf count.

\Cref{sec:comparison} compares the resulting bounds with the
specified numerical tables of \cite{vxxz,advxxz}, augmented by the
same square point \(p=2.258\) and closed under the ordinary
combination operations stated there.  The new bounds strictly improve
this table baseline for \(0.321334<k<1\).  This comparison concerns
recorded upper estimates.  Extensions to additional estimates retain
their stated separator or envelope hypotheses, and strict improvement
is asserted only where the enlarged baseline itself exceeds two.

\subsection{Organization}

\Cref{sec:preliminaries} establishes the tensor and arithmetic-exponent
conventions.  \Cref{alg:section} proves the shared separation lemma
and the finite tensor constructions.  The square argument begins in
\cref{val:section,ap:section}, which define the attainable values and
prove the count-window comparison.
\Cref{cert:w-seed-section} establishes the three-state boundary
estimate used by the numerical certificate in \cref{cert:section}.
\Cref{sec:conclusion} extracts the finite
witness proving the square part of \cref{thm:main}.
\Cref{sec:trees,sec:completion,sec:optimization} develop conditional
entropy, completion trees, and exact finite-family optimization.
\Cref{sec:certificates} proves the remaining parts of
\cref{thm:main}, and \cref{sec:comparison} gives the baseline
comparisons.  Section~\ref{sec:characteristic-transfer} proves the
finite-exception field corollary from two exact rank witnesses and
the characteristic-zero dual corollary using accuracy-dependent witnesses.
The appendices give barrier coefficients and reproduction details.
